\section{Continuous Assessment 2 (CA-2) --- Set D}
\label{sec:ca2_set_d}

\LPUPaperHeader{Continuous Assessment (CA-2)}{D}{50 Minutes}{30 Marks}{CO4 (Multivariable Limits, Continuity, Partial Derivatives \& Optimization), CO5 (Multiple Integrals \& Transformations)}

\begin{center}
\sffamily\textbf{\color{AlertRed}Note: Attempt all six questions. Each question carries 5 marks.}
\end{center}

% ------------------------------------------------------------------------------
% QUESTION PAPER
% ------------------------------------------------------------------------------

\questionitem{1}{
    Expand $f(x,y) = e^x \cos y$ in powers of $x$ and $y$ up to terms of degree 3 using \textbf{Taylor's Theorem for two variables} about the origin $(0,0)$.
}{5}{CO4}{L3 Apply}

\questionitem{2}{
    Find and classify all the \textbf{critical points} (local maxima, local minima, and saddle points) of the function:
    \[
    f(x,y) = x^3 + y^3 - 3xy
    \]
}{5}{CO4}{L3 Apply}

\questionitem{3}{
    If $z = f(x,y)$ where $x = u^2 - v^2$ and $y = 2uv$, prove using the \textbf{multivariable chain rule} that:
    \[
    \left(\frac{\partial z}{\partial u}\right)^2 + \left(\frac{\partial z}{\partial v}\right)^2 = 4(u^2 + v^2)\left[ \left(\frac{\partial z}{\partial x}\right)^2 + \left(\frac{\partial z}{\partial y}\right)^2 \right]
    \]
}{5}{CO4}{L3 Apply}

\questionitem{4}{
    Using \textbf{cylindrical polar coordinates}, evaluate the total volume of the solid enclosed between the paraboloid $z = x^2 + y^2$ and the horizontal plane $z = 4$.
}{5}{CO5}{L3 Apply}

\questionitem{5}{
    Change the order of integration and evaluate:
    \[
    \int_0^a \int_0^{\sqrt{a^2 - y^2}} \sqrt{a^2 - x^2 - y^2} \, dx \, dy
    \]
}{5}{CO5}{L3 Apply}

\questionitem{6}{
    Evaluate the volume of the tetrahedron in the first octant bounded by the coordinate planes and the plane:
    \[
    \frac{x}{a} + \frac{y}{b} + \frac{z}{c} = 1 \quad (a, b, c > 0)
    \]
    using triple integration.
}{5}{CO5}{L3 Apply}

\vspace{5mm}
\hrule
\vspace{5mm}

% ------------------------------------------------------------------------------
% COMPLETE STEP-BY-STEP SOLUTIONS
% ------------------------------------------------------------------------------
\subsection*{Solutions \& Marking Scheme (CA-2 Set D)}

% Solution 1
\begin{solutionbox}[Solution to Question 1 (5 Marks)]
\textbf{Step 1: Taylor's Expansion Formula at $(0,0)$}
\[
f(x,y) = f(0,0) + [x f_x + y f_y] + \frac{1}{2!}[x^2 f_{xx} + 2xy f_{xy} + y^2 f_{yy}] + \frac{1}{3!}[x^3 f_{xxx} + 3x^2 y f_{xxy} + 3x y^2 f_{xyy} + y^3 f_{yyy}] + \dots
\]
\textbf{Step 2: Evaluating Derivatives at $(0,0)$}
For $f(x,y) = e^x \cos y$:
\begin{itemize}[leftmargin=6mm]
    \item $f(0,0) = e^0 \cos 0 = 1$.
    \item $f_x = e^x \cos y \implies f_x(0,0) = 1$.
    \item $f_y = -e^x \sin y \implies f_y(0,0) = 0$.
    \item $f_{xx} = e^x \cos y \implies f_{xx}(0,0) = 1$.
    \item $f_{xy} = -e^x \sin y \implies f_{xy}(0,0) = 0$.
    \item $f_{yy} = -e^x \cos y \implies f_{yy}(0,0) = -1$.
    \item $f_{xxx} = e^x \cos y \implies f_{xxx}(0,0) = 1$.
    \item $f_{xxy} = -e^x \sin y \implies f_{xxy}(0,0) = 0$.
    \item $f_{xyy} = -e^x \cos y \implies f_{xyy}(0,0) = -1$.
    \item $f_{yyy} = e^x \sin y \implies f_{yyy}(0,0) = 0$.
\end{itemize}
\textbf{Step 3: Assembling the Series}
\[
f(x,y) = 1 + x + \frac{1}{2}(x^2 - y^2) + \frac{1}{6}(x^3 - 3x y^2) + \dots
\]
\[
\mathbf{f(x,y) = 1 + x + \frac{x^2}{2} - \frac{y^2}{2} + \frac{x^3}{6} - \frac{x y^2}{2} + \dots}
\]

\begin{markingrubric}
\textbf{Mark Distribution:}
$\bullet$ Statement of Taylor's two-variable expansion formula: \textbf{1.5 Marks}\\
$\bullet$ Correct evaluation of all partial derivatives up to order 3 at $(0,0)$: \textbf{2 Marks}\\
$\bullet$ Assembled series expression: \textbf{1.5 Marks}
\end{markingrubric}
\end{solutionbox}

% Solution 2
\begin{solutionbox}[Solution to Question 2 (5 Marks)]
\textbf{Step 1: First Partial Derivatives}
\[
f_x = 3x^2 - 3y = 0 \implies y = x^2 \label{eq:ext1}
\]
\[
f_y = 3y^2 - 3x = 0 \implies x = y^2 \label{eq:ext2}
\]
Substitute \eqref{eq:ext1} into \eqref{eq:ext2}:
\[
x = (x^2)^2 = x^4 \implies x(x^3 - 1) = 0 \implies x = 0 \quad \text{or} \quad x = 1
\]
For $x = 0 \implies y = 0^2 = 0 \implies \mathbf{P_1(0,0)}$.\\
For $x = 1 \implies y = 1^2 = 1 \implies \mathbf{P_2(1,1)}$.

\textbf{Step 2: Second Partial Derivatives \& Discriminant}
\[
r = f_{xx} = 6x, \quad s = f_{xy} = -3, \quad t = f_{yy} = 6y
\]
\[
D = rt - s^2 = (6x)(6y) - (-3)^2 = 36xy - 9
\]

\textbf{Step 3: Classification}
\begin{itemize}[leftmargin=6mm]
    \item At point $P_1(0,0)$:
    \[
    D = 36(0)(0) - 9 = -9 < 0
    \]
    Since $D < 0$, $\mathbf{(0,0)}$ is a \textbf{Saddle Point}.
    \item At point $P_2(1,1)$:
    \[
    D = 36(1)(1) - 9 = 27 > 0 \quad \text{and} \quad r = 6(1) = 6 > 0
    \]
    Since $D > 0$ and $r > 0$, $\mathbf{(1,1)}$ is a point of \textbf{Local Minimum}.
    The local minimum value is:
    \[
    f(1,1) = 1^3 + 1^3 - 3(1)(1) = 1 + 1 - 3 = \mathbf{-1}
    \]
\end{itemize}

\begin{markingrubric}
\textbf{Mark Distribution:}
$\bullet$ First partial derivatives and finding points $(0,0)$ and $(1,1)$: \textbf{2 Marks}\\
$\bullet$ Computing $r, s, t$ and discriminant formula $D = rt - s^2$: \textbf{1.5 Marks}\\
$\bullet$ Classifying $(0,0)$ as saddle and $(1,1)$ as minimum with value $-1$: \textbf{1.5 Marks}
\end{markingrubric}
\end{solutionbox}

% Solution 3
\begin{solutionbox}[Solution to Question 3 (5 Marks)]
\textbf{Step 1: Chain Rule Expressions}
By the multivariable chain rule:
\[
\frac{\partial z}{\partial u} = \frac{\partial z}{\partial x} \frac{\partial x}{\partial u} + \frac{\partial z}{\partial y} \frac{\partial y}{\partial u}
\]
\[
\frac{\partial z}{\partial v} = \frac{\partial z}{\partial x} \frac{\partial x}{\partial v} + \frac{\partial z}{\partial y} \frac{\partial y}{\partial v}
\]
Given $x = u^2 - v^2$ and $y = 2uv$:
\[
\frac{\partial x}{\partial u} = 2u, \quad \frac{\partial x}{\partial v} = -2v, \quad \frac{\partial y}{\partial u} = 2v, \quad \frac{\partial y}{\partial v} = 2u
\]
Substitute these into the chain rule:
\[
\frac{\partial z}{\partial u} = 2u \frac{\partial z}{\partial x} + 2v \frac{\partial z}{\partial y}
\]
\[
\frac{\partial z}{\partial v} = -2v \frac{\partial z}{\partial x} + 2u \frac{\partial z}{\partial y}
\]
\textbf{Step 2: Squaring and Adding}
\[
\left(\frac{\partial z}{\partial u}\right)^2 = 4\left[ u^2 \left(\frac{\partial z}{\partial x}\right)^2 + v^2 \left(\frac{\partial z}{\partial y}\right)^2 + 2uv \frac{\partial z}{\partial x}\frac{\partial z}{\partial y} \right]
\]
\[
\left(\frac{\partial z}{\partial v}\right)^2 = 4\left[ v^2 \left(\frac{\partial z}{\partial x}\right)^2 + u^2 \left(\frac{\partial z}{\partial y}\right)^2 - 2uv \frac{\partial z}{\partial x}\frac{\partial z}{\partial y} \right]
\]
Adding the two equations, the cross terms cancel:
\[
\left(\frac{\partial z}{\partial u}\right)^2 + \left(\frac{\partial z}{\partial v}\right)^2 = 4\left[ (u^2 + v^2)\left(\frac{\partial z}{\partial x}\right)^2 + (u^2 + v^2)\left(\frac{\partial z}{\partial y}\right)^2 \right]
\]
Factor out $(u^2 + v^2)$:
\[
\mathbf{\left(\frac{\partial z}{\partial u}\right)^2 + \left(\frac{\partial z}{\partial v}\right)^2 = 4(u^2 + v^2)\left[ \left(\frac{\partial z}{\partial x}\right)^2 + \left(\frac{\partial z}{\partial y}\right)^2 \right]} \quad \text{(Hence Proved!)}
\]

\begin{markingrubric}
\textbf{Mark Distribution:}
$\bullet$ Stating chain rule and partial derivatives of $x$ and $y$: \textbf{2 Marks}\\
$\bullet$ Setting up algebraic expressions for squared derivatives: \textbf{1.5 Marks}\\
$\bullet$ Cancellation of cross-product terms and factorization: \textbf{1.5 Marks}
\end{markingrubric}
\end{solutionbox}

% Solution 4
\begin{solutionbox}[Solution to Question 4 (5 Marks)]
\textbf{Step 1: Cylindrical Coordinates Formulation}
In cylindrical coordinates: $x = r\cos\theta, y = r\sin\theta, z = z$, with $dV = r \, dz \, dr \, d\theta$.
The bounding surfaces are:
\begin{itemize}[leftmargin=6mm]
    \item Bottom surface: Paraboloid $z = r^2$.
    \item Top surface: Horizontal plane $z = 4$.
\end{itemize}
Intersection of the surfaces occurs when $r^2 = 4 \implies r = 2$.
\textbf{Step 2: Integration Limits}
\[
\theta \in [0, 2\pi], \quad r \in [0, 2], \quad z \in [r^2, 4]
\]
\textbf{Step 3: Volume Evaluation}
\[
V = \int_0^{2\pi} \int_0^2 \int_{r^2}^4 r \, dz \, dr \, d\theta = \left(\int_0^{2\pi} d\theta\right) \int_0^2 r [z]_{r^2}^4 \, dr
\]
\[
= 2\pi \int_0^2 r (4 - r^2) \, dr = 2\pi \int_0^2 (4r - r^3) \, dr = 2\pi \left[ 2r^2 - \frac{r^4}{4} \right]_0^2
\]
\[
= 2\pi \left[ 2(4) - \frac{16}{4} \right] = 2\pi [8 - 4] = 2\pi(4) = \mathbf{8\pi}
\]

\begin{markingrubric}
\textbf{Mark Distribution:}
$\bullet$ Setting up cylindrical coordinates and intersection radius $r = 2$: \textbf{1.5 Marks}\\
$\bullet$ Correct bounds for $z, r, \theta$: \textbf{1.5 Marks}\\
$\bullet$ Integration yielding $8\pi$: \textbf{2 Marks}
\end{markingrubric}
\end{solutionbox}

% Solution 5
\begin{solutionbox}[Solution to Question 5 (5 Marks)]
\textbf{Step 1: Geometry of the Region}
Original limits:
\[
y \in [0, a], \quad x \in [0, \sqrt{a^2 - y^2}]
\]
Since $x^2 = a^2 - y^2 \implies x^2 + y^2 = a^2$, this is the first quadrant of the circle of radius $a$.
Converting directly to polar coordinates:
\[
x = r\cos\theta, \quad y = r\sin\theta, \quad dx \, dy = r \, dr \, d\theta
\]
\[
r \in [0, a], \quad \theta \in \left[0, \frac{\pi}{2}\right]
\]
\textbf{Step 2: Transforming the Integrand}
$\sqrt{a^2 - x^2 - y^2} = \sqrt{a^2 - r^2}$.
\textbf{Step 3: Evaluating the Integral}
\[
I = \int_0^{\pi/2} d\theta \cdot \int_0^a \sqrt{a^2 - r^2} \cdot r \, dr
\]
\[
\int_0^{\pi/2} d\theta = \frac{\pi}{2}
\]
For the radial integral, let $u = a^2 - r^2 \implies du = -2r \, dr \implies r \, dr = -\frac{du}{2}$:
\[
\int_0^a r\sqrt{a^2 - r^2} \, dr = \int_{a^2}^0 u^{1/2} \left(-\frac{du}{2}\right) = \frac{1}{2} \int_0^{a^2} u^{1/2} \, du = \frac{1}{2} \left[ \frac{u^{3/2}}{3/2} \right]_0^{a^2} = \frac{1}{3} a^3
\]
Multiply the results:
\[
I = \left(\frac{\pi}{2}\right) \left(\frac{1}{3} a^3\right) = \mathbf{\frac{\pi a^3}{6}}
\]

\begin{markingrubric}
\textbf{Mark Distribution:}
$\bullet$ Identifying circular region in first quadrant: \textbf{1.5 Marks}\\
$\bullet$ Polar coordinate conversion with Jacobian: \textbf{1.5 Marks}\\
$\bullet$ Integration steps yielding $\pi a^3/6$: \textbf{2 Marks}
\end{markingrubric}
\end{solutionbox}

% Solution 6
\begin{solutionbox}[Solution to Question 6 (5 Marks)]
\textbf{Step 1: Region of the Tetrahedron}
Bounded by $x \ge 0, y \ge 0, z \ge 0$ and $\frac{x}{a} + \frac{y}{b} + \frac{z}{c} \le 1$.
The limits of integration are:
\[
z \in \left[0, c\left(1 - \frac{x}{a} - \frac{y}{b}\right)\right]
\]
\[
y \in \left[0, b\left(1 - \frac{x}{a}\right)\right]
\]
\[
x \in [0, a]
\]
\textbf{Step 2: Integration with respect to $z$}
\[
\int_0^{c(1 - x/a - y/b)} dz = c\left(1 - \frac{x}{a} - \frac{y}{b}\right)
\]
\textbf{Step 3: Integration with respect to $y$}
\[
c \int_0^{b(1 - x/a)} \left(1 - \frac{x}{a} - \frac{y}{b}\right) dy = c \left[ \left(1 - \frac{x}{a}\right)y - \frac{y^2}{2b} \right]_0^{b(1 - x/a)}
\]
\[
= c \left[ b\left(1 - \frac{x}{a}\right)^2 - \frac{b^2(1 - x/a)^2}{2b} \right] = \frac{bc}{2}\left(1 - \frac{x}{a}\right)^2
\]
\textbf{Step 4: Integration with respect to $x$}
\[
V = \frac{bc}{2} \int_0^a \left(1 - \frac{x}{a}\right)^2 dx = \frac{bc}{2} \left[ -\frac{a}{3}\left(1 - \frac{x}{a}\right)^3 \right]_0^a = \frac{bc}{2}\left(0 - \left(-\frac{a}{3}\right)\right) = \mathbf{\frac{abc}{6}}
\]

\begin{markingrubric}
\textbf{Mark Distribution:}
$\bullet$ Correct bounds for $x, y, z$: \textbf{1.5 Marks}\\
$\bullet$ Sequential integration with respect to $z$ and $y$: \textbf{2 Marks}\\
$\bullet$ Final integral evaluation giving $V = abc/6$: \textbf{1.5 Marks}
\end{markingrubric}
\end{solutionbox}

\newpage
